India's Mars Orbiter Mission (MOM) was launched using the Polar Satellite
Launch Vehicle (PSLV) on 5 November 2013. The dry mass of MOM (body +
instruments) was 500\,kg and it carried fuel of mass 852\,kg. It was initially
placed in an elliptical orbit around the Earth with perigee at a height of
264.1\,km and apogee at a height of 23903.6\,km, above the surface of the
Earth. After raising the orbit six times, MOM was transferred to a trans-Mars
injection orbit (Hohmann orbit).

The first such orbit-raising was performed by firing the engines for a very
short time near the perigee. The engines were fired to change the orbit
without changing the plane of the orbit and without changing its perigee. This
gave a net impulse of $1.73 \times 10^5$\,kg\,m\,s$^{-1}$ to the satellite.
Ignore the change in mass due to burning of fuel.

\begin{description}
\item[(T9.1)] What is the height of the new apogee, $h_{\mathrm{a}}$ above the
  surface of the Earth, after this engine burn? \hfill[14]
\item[(T9.2)] Find the eccentricity ($e$) of the new orbit after the burn and
  the new orbital period ($P$) of MOM in hours. \hfill[6]
\end{description}
